\renewcommand{\thetheorem}{1.\arabic{theorem}}
\setcounter{theorem}{0}
\setcounter{chapter}{0}
% ---------------------------------------------------------------------------
% Provenance of the prose in this file.
%   % Extractor: <model>   the block below was transcribed from Prof. Biran's
%                          handwritten notes by that model
%   % Creator:   <model>   the block below was written by that model and has no
%                          counterpart in the notes
% A tag holds until the next tag.
%
% This chapter was written before the convention existed, so the prose carries
% no tags except on the blocks added by later passes. Untagged prose is the
% original transcription.
% ---------------------------------------------------------------------------

\chapter{The Fibonacci Sequence}
\stepcounter{section}
\subsection{Introduction and Generalization}

We begin our study with the classical Fibonacci sequence, a sequence of numbers that naturally emerges in various mathematical and biological contexts.

\begin{definition}[The Fibonacci Sequence]
\label{def:classical_fibonacci}
The \newterm{Fibonacci sequence} is an infinite sequence of integers $a_0, a_1, a_2, a_3, \dots, a_n, \dots$ defined by the following recurrence relation:
\begin{equation} \label{eq:standard_fib_recurrence}
\begin{cases}
a_0 := 0 \\
a_1 := 1 \\
a_n := a_{n-1} + a_{n-2} & \text{for } n \ge 2
\end{cases}
\end{equation}
\end{definition}

\begin{remark}
This sequence was introduced to Western mathematics by Leonardo Fibonacci, a medieval mathematician from Pisa (1170--1250), to model the growth of rabbit populations. Computing the first few terms, we easily find:
\[
a_0 = 0, \quad a_1 = 1, \quad a_2 = 1, \quad a_3 = 2, \quad a_4 = 3, \quad a_5 = 5, \quad a_6 = 8, \quad a_7 = 13, \quad a_8 = 21, \dots
\]
\end{remark}

A natural, central question arises: Is there a closed, explicit formula for the general term $a_n$? To answer this, let us briefly look at other well-known mathematical sequences for comparison:
\begin{enumerate}
    \item \textbf{Arithmetic Sequence:} Defined by $a_0 := a$ and $a_n := a_{n-1} + d$ for $n \ge 1$. The explicit formula for the general term is given by $a_n = a + n \cdot d$.
    \item \textbf{Geometric Sequence:} Defined by $a_0 := a$ and $a_n := a_{n-1} \cdot q$ for $n \ge 1$. The explicit formula is given by $a_n = a \cdot q^n$.
\end{enumerate}

Clearly, the recurrence relation for the Fibonacci sequence is fundamentally different. To find a formula for $a_n$, it is highly beneficial to step back and generalize the problem.

\begin{definition}[Generalized Fibonacci Sequence]
\label{def:generalized_fibonacci}
Let $a_0, a_1 \in \mathbb{R}$ be given initial real numbers. We define a sequence $\mathcal{A}$ to be a \newterm{Fibonacci sequence} if it satisfies the relation:
\begin{equation} \label{eq:generalized_fib_recurrence}
\begin{cases}
\alpha_0 := a_0 \\
\alpha_1 := a_1 \\
\alpha_n := \alpha_{n-1} + \alpha_{n-2} & \text{for } n \ge 2
\end{cases}
\end{equation}
We denote such a sequence uniquely determined by its starting values as $\mathcal{F}_{a_0, a_1}$. Furthermore, we denote the collection (or space) of all such generalized Fibonacci sequences by $\Fib$.
\end{definition}

Notice that our original, classical Fibonacci sequence is precisely $\mathcal{F}_{0, 1}$. 

\subsection{The Algebraic Structure of \texorpdfstring{$\Fib$}{Fib}}

The collection $\Fib$ possesses a very rich internal algebraic structure, which we will exploit to find our explicit formula.

\begin{proposition}[Closure Properties of $\Fib$] \label{prop:fib_closure}
The collection of Fibonacci sequences, $\Fib$, is closed under term-wise addition and scalar multiplication. Specifically,
\begin{enumerate}
    \item If $\mathcal{A}_1, \mathcal{A}_2 \in \Fib$, then $\mathcal{A}_1 + \mathcal{A}_2 \in \Fib$.
    \item If $\mathcal{A} \in \Fib$ and $\alpha \in \mathbb{R}$, then $\alpha \mathcal{A} \in \Fib$.
\end{enumerate}
\end{proposition}

\begin{proof}
We address each closure property:
\begin{enumerate}[label=\textbf{(\alph*)}]
    \item Let $\mathcal{A}_1 = (a_0, a_1, a_2, \dots)$ and $\mathcal{A}_2 = (b_0, b_1, b_2, \dots)$ be two sequences in $\Fib$. We define their sum term-wise as $\mathcal{A}_1 + \mathcal{A}_2 := (c_0, c_1, c_2, \dots)$, where $c_n := a_n + b_n$. We must verify that the sequence $(c_n)$ obeys the Fibonacci recurrence relation, namely $c_n = c_{n-1} + c_{n-2}$ for $n \ge 2$. By algebraic manipulation, we find:
\[
c_{n-1} + c_{n-2} = (a_{n-1} + b_{n-1}) + (a_{n-2} + b_{n-2}) = (a_{n-1} + a_{n-2}) + (b_{n-1} + b_{n-2})
\]
Because both $\mathcal{A}_1$ and $\mathcal{A}_2$ are Fibonacci sequences, we can substitute $a_{n-1} + a_{n-2} = a_n$ and $b_{n-1} + b_{n-2} = b_n$. Therefore,
\[
c_{n-1} + c_{n-2} = a_n + b_n = c_n
\]
This shows that $\mathcal{F}_{a_0, a_1} + \mathcal{F}_{b_0, b_1} = \mathcal{F}_{a_0 + b_0, a_1 + b_1}$, and therefore, $\Fib$ is closed under addition.

    \item Let $\mathcal{A} = (a_0, a_1, a_2, \dots) \in \Fib$ and $\alpha \in \mathbb{R}$. We define $\alpha \mathcal{A} := (\alpha a_0, \alpha a_1, \alpha a_2, \dots)$. We need to check that $\alpha a_n = \alpha a_{n-1} + \alpha a_{n-2}$. This is true because $a_n = a_{n-1} + a_{n-2}$, and we simply multiply both sides by $\alpha$. In fact, we see that $\alpha \mathcal{F}_{a_0, a_1} = \mathcal{F}_{\alpha a_0, \alpha a_1}$, so $\Fib$ is closed under scalar multiplication.
\end{enumerate}
\end{proof}

\begin{corollary}
\label{cor:fib_linear_combination}
By \cref{prop:fib_closure}, linear combinations of Fibonacci sequences are also Fibonacci sequences. That is, for any scalars $\alpha, \beta \in \mathbb{R}$, we have:
\[
\alpha \mathcal{F}_{a_0, a_1} + \beta \mathcal{F}_{b_0, b_1} = \mathcal{F}_{\alpha a_0 + \beta b_0, \alpha a_1 + \beta b_1}
\]
\end{corollary}

\begin{remark}
As a direct consequence of this structure, in order to find a general formula for \textit{any} sequence $\mathcal{F}_{a_0, a_1}$, it is enough to find explicit formulas for the two foundational sequences $\mathcal{F}_{1,0}$ and $\mathcal{F}_{0,1}$. This is because any sequence can be decomposed as:
\[
\mathcal{F}_{a_0, a_1} = a_0 \cdot \mathcal{F}_{1,0} + a_1 \cdot \mathcal{F}_{0,1}
\]
We will see later in the course that $\Fib$ represents a $2$-dimensional vector space. Every element in this space can be expressed by a linear combination of some two base elements (but $1$ element alone is not enough). What precisely constitutes a \qt{vector space} and \qt{dimension} will be formally defined later in the course. The corresponding picture is drawn in \cref{fig:fib_plane}.
\end{remark}

\begin{exercise}[Distinguishing Fibonacci Sequences by Random Samples]
\label{exc:fib_dimension_two_samples}
Let $F = (\alpha_0, \alpha_1, \alpha_2, \dots)$ and $G = (\beta_0, \beta_1, \beta_2, \dots)$ be two sequences in $\Fib$. The sequences themselves are not given explicitly, but for any natural number $i$, you can test whether $\alpha_i = \beta_i$. You can query as many distinct indices as you wish.

\textbf{Claim:} It is sufficient to check equality at two distinct indices $i \ne j$ to determine with certainty whether $F = G$. Explain why this claim is true.
\exinfo{This exercise is Problem 6 of Problem Sheet 1, Linear Algebra I, Autumn Semester 2025.}
\exsol{sol:fib_dimension_two_samples}
\end{exercise}

\subsection{Deriving the Explicit Formula}

Perhaps one of the sequences in $\Fib$ matches a form we already know? 

First, we can rule out arithmetic sequences. Except for the trivial zero sequence $\mathcal{F}_{0,0}$, no Fibonacci sequence can be arithmetic. The reason is straightforward: for an arithmetic sequence, the difference between consecutive terms must be constant, meaning $a_n - a_{n-1} = d$. However, for a Fibonacci sequence, $a_n - a_{n-1} = a_{n-2}$. This difference grows and cannot remain constant. Thus, arithmetic sequences are unsuitable for this purpose.
% Creator: Opus 5
(Prof. Biran leaves the details of this as \cref{exc:ch1_list} \textbf{(d)}, and they are worth having, because the appeal to growth is not quite the argument: what the two displays give is that $a_{n-2} = d$ for every $n \ge 2$, that is, that the sequence is \emph{constant}, equal to $d$ throughout. A constant sequence obeys the Fibonacci recurrence exactly when $d = d + d$, which forces $d = 0$. No estimate of growth is needed, and the conclusion holds for negative and non-monotone sequences alike).

Next, we ask: perhaps for some choice of $a_0, a_1$, the sequence $\mathcal{F}_{a_0, a_1}$ coincides with a geometric sequence? Let us try a geometric sequence of the form $\mathcal{G} = (1, q, q^2, q^3, \dots, q^n, \dots)$. 

\begin{lemma}
\label{lem:geometric_fibonacci}
A geometric sequence $\mathcal{G} = (1, q, q^2, \dots)$ with $q \ne 0$ is a Fibonacci sequence if and only if $q = \frac{1 \pm \sqrt{5}}{2}$.
\end{lemma}

\begin{proof}
For $\mathcal{G}$ to be a Fibonacci sequence, it must satisfy the recurrence relation for all $n \ge 2$:
\[
q^n = q^{n-1} + q^{n-2}
\]
Since $q \ne 0$, we can divide the entire equation by $q^{n-2}$, which simplifies the condition to a quadratic equation:
\[
q^2 = q + 1 \implies q^2 - q - 1 = 0
\]
Applying the standard quadratic formula yields two possible solutions:
\[
q = \frac{1 \pm \sqrt{5}}{2}
\]
We denote these two specific roots as $\varphi := \frac{1 + \sqrt{5}}{2} \approx 1.618$ (the golden ratio) and $\psi := \frac{1 - \sqrt{5}}{2} \approx -0.618$. Consequently, we have discovered two purely geometric sequences that belong to $\Fib$:
\[
\mathcal{F}_{1, \varphi} = (1, \varphi, \varphi^2, \varphi^3, \dots, \varphi^n, \dots)
\]
\[
\mathcal{F}_{1, \psi} = (1, \psi, \psi^2, \psi^3, \dots, \psi^n, \dots)
\]
\end{proof}

\begin{figure}[ht!]
  \centering
  \begin{tikzpicture}[scale=1.35, >=Stealth]
    %, background grid ---
    \draw[step=0.5, gray!15, very thin] (-1.7,-1.2) rectangle (2.2,2.2);
    %, axes ---
    \draw[->, gray!70] (-1.80,0) -- (2.45,0) node[right, black, font=\small] {$a_0$};
    \draw[->, gray!70] (0,-1.30) -- (0,2.40) node[above, black, font=\small] {$a_1$};
    %, the two eigen-directions (drawn faintly here, revisited later) ---
    \draw[ThemeAccent!55, thick] (-0.74,-1.20) -- (1.36,2.20);
    \node[ThemeAccent, above left, font=\scriptsize] at (1.30,2.06) {ratio $\varphi$};
    \draw[ThemeGreen!55, thick] (-1.70,1.05) -- (1.94,-1.20);
    \node[ThemeGreen, below left, font=\scriptsize] at (1.92,-1.06) {ratio $\psi$};
    %, the two building blocks ---
    \draw[->, ThemePurple, very thick] (0,0) -- (1,0);
    \fill[ThemePurple] (1,0) circle (1.7pt);
    \node[ThemePurple, below, font=\scriptsize] at (1.02,-0.06) {$\mathcal{F}_{1,0}$};
    \draw[->, ThemePurple, very thick] (0,0) -- (0,1);
    \fill[ThemePurple] (0,1) circle (1.7pt);
    \node[ThemePurple, left, font=\scriptsize] at (-0.06,1.02) {$\mathcal{F}_{0,1}$};
    %, the two geometric Fibonacci sequences ---
    \fill[ThemeAccent] (1,1.618) circle (1.7pt);
    \node[ThemeAccent, right, font=\scriptsize] at (1.06,1.64) {$\mathcal{F}_{1,\varphi}$};
    \fill[ThemeGreen] (1,-0.618) circle (1.7pt);
    \node[ThemeGreen, right, font=\scriptsize] at (1.06,-0.64) {$\mathcal{F}_{1,\psi}$};
    %, a general element, decomposed along the two building blocks ---
    \fill[EmphBlue] (1.7,0.8) circle (1.7pt);
    \node[EmphBlue, above right, font=\scriptsize] at (1.66,0.82) {$\mathcal{F}_{a_0,a_1}$};
    \draw[EmphBlue, dotted, thick] (1.7,0) -- (1.7,0.8) -- (0,0.8);
    \node[EmphBlue, below, font=\scriptsize] at (1.7,-0.04) {$a_0$};
    \node[EmphBlue, left, font=\scriptsize] at (-0.04,0.8) {$a_1$};
    \fill (0,0) circle (1.4pt);
    \node[below left, font=\scriptsize] at (0,0) {$\mathcal{F}_{0,0}$};
  \end{tikzpicture}
  \caption{The space $\Fib$ drawn as a plane: a Fibonacci sequence is completely
  determined by its two initial values, so the point $(a_0, a_1)$ \textit{is} the
  sequence $\mathcal{F}_{a_0,a_1}$. Addition and rescaling of sequences become the
  usual addition and rescaling of points, which is the content of
  \cref{prop:fib_closure}, and the decomposition
  $\mathcal{F}_{a_0,a_1} = a_0 \mathcal{F}_{1,0} + a_1 \mathcal{F}_{0,1}$ is just
  reading off coordinates. The two coloured lines carry the geometric Fibonacci
  sequences found in \cref{lem:geometric_fibonacci}; they will turn out to be the
  natural axes of this plane.}
  \label{fig:fib_plane}
\end{figure}

\begin{remark}
The number $\varphi$ is one of the most persistent constants in mathematics. The
defining relation $\varphi^2 = \varphi + 1$ can be rewritten as
$\varphi = 1 + \tfrac{1}{\varphi}$, which says that removing a unit square from a
$\varphi \times 1$ rectangle leaves a rectangle of exactly the same proportions.
Iterating this observation produces the classical subdivision of a rectangle into
squares whose side lengths are consecutive Fibonacci numbers, drawn in
\cref{fig:golden_spiral}. The quarter circles inscribed into those squares fit
together smoothly, and the resulting spiral grows by a factor of $\varphi$ per quarter
turn.
\end{remark}

\begin{figure}[ht!]
  \centering
  \begin{tikzpicture}[scale=0.50]
    %, the Fibonacci squares tiling a 13 x 8 rectangle ---
    \draw[ThemePurple!55, fill=ThemePurple!5]  (5,0) rectangle (13,8);
    \draw[ThemePurple!55, fill=ThemePurple!9]  (0,0) rectangle (5,5);
    \draw[ThemePurple!55, fill=ThemePurple!13] (0,5) rectangle (3,8);
    \draw[ThemePurple!55, fill=ThemePurple!17] (3,6) rectangle (5,8);
    \draw[ThemePurple!55, fill=ThemePurple!21] (3,5) rectangle (4,6);
    \draw[ThemePurple!55, fill=ThemePurple!25] (4,5) rectangle (5,6);
    %, side lengths ---
    \node[ThemePurple, font=\small] at (9,4)     {$8$};
    \node[ThemePurple, font=\small] at (2.5,2.5) {$5$};
    \node[ThemePurple, font=\small] at (1.5,6.5) {$3$};
    \node[ThemePurple, font=\small] at (4,7)     {$2$};
    \node[ThemePurple, font=\tiny]  at (3.5,5.5) {$1$};
    \node[ThemePurple, font=\tiny]  at (4.5,5.5) {$1$};
    %, the golden spiral: five quarter turns, each of radius a Fibonacci number ---
    \draw[ThemeAccent, very thick] (3,6) arc[start angle=180, end angle=360, radius=1];
    \draw[ThemeAccent, very thick] (5,6) arc[start angle=0,   end angle=90,  radius=2];
    \draw[ThemeAccent, very thick] (3,8) arc[start angle=90,  end angle=180, radius=3];
    \draw[ThemeAccent, very thick] (0,5) arc[start angle=180, end angle=270, radius=5];
    \draw[ThemeAccent, very thick] (5,0) arc[start angle=270, end angle=360, radius=8];
    %, the enclosing 13 x 8 rectangle ---
    \draw[ThemePurple, thick] (0,0) rectangle (13,8);
    \draw[<->, gray!80] (0,-0.9) -- (13,-0.9)
      node[midway, text=black, fill=white, inner sep=1pt, font=\scriptsize] {$13 = a_7$};
    \draw[<->, gray!80] (13.9,0) -- (13.9,8)
      node[midway, text=black, fill=white, inner sep=1pt, font=\scriptsize] {$8 = a_6$};
  \end{tikzpicture}
  \caption{Squares with consecutive Fibonacci side lengths $1, 1, 2, 3, 5, 8$ tile a
  $13 \times 8$ rectangle, because $a_{n-1} + a_n = a_{n+1}$ is exactly the statement
  that a square of side $a_n$ glued onto an $a_{n-1} \times a_n$ rectangle produces an
  $a_n \times a_{n+1}$ rectangle. Since $a_{n+1}/a_n \to \varphi$, these rectangles
  become ever closer to the golden rectangle.}
  \label{fig:golden_spiral}
\end{figure}

\begin{theorem}[Binet's Formula]
\label{thm:binet_formula}
The explicit formula for the $n$-th term of the classical Fibonacci sequence $\mathcal{F}_{0,1}$ is given by \newterm{Binet's Formula}:
\[
a_n = \frac{1}{\sqrt{5}} \left( \left(\frac{1+\sqrt{5}}{2}\right)^n - \left(\frac{1-\sqrt{5}}{2}\right)^n \right)
\]
\end{theorem}

\begin{proof}
So far, by \cref{lem:geometric_fibonacci}, we have found formulas for two specific sequences, $\mathcal{F}_{1, \varphi}$ and $\mathcal{F}_{1, \psi}$. Our original goal is to find a formula for the classical sequence $\mathcal{F}_{0,1}$.

We use the linearity principle established in \cref{cor:fib_linear_combination}. We seek scalars $\alpha, \beta \in \mathbb{R}$ such that:
\[
\alpha \cdot \mathcal{F}_{1, \varphi} + \beta \cdot \mathcal{F}_{1, \psi} = \mathcal{F}_{0,1}
\]
Comparing the first two terms ($n=0$ and $n=1$) of these sequences yields a system of linear equations:
\begin{equation}
\begin{cases}
\alpha \cdot 1 + \beta \cdot 1 = 0 \\
\alpha \cdot \varphi + \beta \cdot \psi = 1
\end{cases}
\end{equation}
From the first equation, we immediately deduce that $\beta = -\alpha$. Substituting this into the second equation gives:
\[
\alpha \cdot \varphi - \alpha \cdot \psi = 1 \implies \alpha (\varphi - \psi) = 1
\]
Notice that the difference $\varphi - \psi = \sqrt{5}$. Therefore, $\alpha = 1/\sqrt{5}$ and $\beta = -1/\sqrt{5}$. This leads directly to:
\[
a_n = \frac{1}{\sqrt{5}} (\varphi^n - \psi^n) = \frac{1}{\sqrt{5}} \left( \left(\frac{1+\sqrt{5}}{2}\right)^n - \left(\frac{1-\sqrt{5}}{2}\right)^n \right)
\]
\end{proof}

\begin{exercise}[The Golden Ratio and Generalized Fibonacci Sequences]
\label{exc:golden_ratio_part1}
Recall that the classical Fibonacci sequence $\mathcal{F}_{0,1} = (a_0, a_1, \dots)$ satisfies the recurrence $a_n = a_{n-1} + a_{n-2}$ with initial values $a_0 = 0$ and $a_1 = 1$.
\begin{enumerate}
    \item Find an explicit closed formula for the sequence $\mathcal{F}_{1,0} = (b_0, b_1, \dots)$ with initial values $b_0 = 1$ and $b_1 = 0$.
    \item Write an explicit formula for the general Fibonacci sequence $\mathcal{F}_{a,b}$ that depends only on the initial values $a, b$ and the index $n$.
\end{enumerate}
\exinfo{This exercise is Problem 1 of Problem Sheet 1, Linear Algebra I, Autumn Semester 2025.}
\exsol{sol:golden_ratio_part1}
\end{exercise}

\subsection{Where the Idea Came From: Symmetries of \texorpdfstring{$\Fib$}{Fib}}

The derivation of \cref{thm:binet_formula} is complete, but it rests on a step that
looks like a lucky guess: why should one ever think of hunting for a \emph{geometric}
sequence inside $\Fib$? The answer is not luck, and uncovering it introduces the single
most important theme of this course.

Let us look at maps
\[
T \colon \Fib \longrightarrow \Fib,
\]
also called functions, transformations or operators. Such a $T$ takes as input an
arbitrary Fibonacci sequence $\mathcal{A}$ and returns as output another Fibonacci
sequence $T(\mathcal{A})$. We have just seen in \cref{prop:fib_closure} that $\Fib$
carries an algebraic structure: it is closed under addition and under multiplication
by scalars. It is therefore natural to restrict attention to those maps that
\emph{respect} this structure.

\begin{definition}[Linear Map on $\Fib$]
\label{def:linear_map_fib}
A map $T \colon \Fib \to \Fib$ is called a \newterm{linear map} if
\begin{equation} \label{eq:fib_linearity}
\begin{cases}
T(\mathcal{A} + \mathcal{B}) = T(\mathcal{A}) + T(\mathcal{B}) & \text{for all } \mathcal{A}, \mathcal{B} \in \Fib, \\[0.4ex]
T(\alpha \mathcal{A}) = \alpha \, T(\mathcal{A}) & \text{for all } \alpha \in \mathbb{R}, \ \mathcal{A} \in \Fib.
\end{cases}
\end{equation}
\end{definition}

There are many examples of such maps, though the first ones that come to mind are not
very exciting:
\begin{enumerate}
    \item the identity $T := \id$, that is $T(\mathcal{A}) := \mathcal{A}$;
    \item the rescaling $T(\mathcal{A}) := \gamma \cdot \mathcal{A}$ for a fixed $\gamma \in \mathbb{R}$.
\end{enumerate}
Neither of these takes any notice of the fact that the elements of $\Fib$ are
\textit{sequences}. Here is another one, which does:

\begin{definition}[The Shift Map]
\label{def:shift_map}
The \newterm{shift map} $S \colon \Fib \to \Fib$ is defined by
\[
S(a_0, a_1, a_2, \dots, a_n, \dots) := (a_1, a_2, \dots, a_n, \dots),
\]
that is, it simply forgets the first term of the sequence.
\end{definition}

\begin{exercise}[Linearity of the Shift Operator]
\label{exc:shift_is_linear}
Check that $S$ is linear.
\exinfo{This exercise is taken from Prof.~Biran's handwritten lecture notes.}
\exsol{sol:shift_is_linear_b}
\end{exercise}

Mathematicians like to probe a map $T \colon X \to X$ by looking for its special, or
extreme, values. Two such notions will accompany us for the rest of the course.

\begin{definition}[Fixed Point and Eigenvector]
\label{def:fixed_point_eigenvector}
Let $T \colon X \to X$ be a map. An element $x_0 \in X$ is a \newterm{fixed point} of
$T$ if $T(x_0) = x_0$. An element $y_0 \in X$ is an \newterm{eigenvector} of $T$ if
$T(y_0)$ is proportional to $y_0$, that is, if
\[
T(y_0) = \lambda \cdot y_0 \quad \text{for some } \lambda \in \mathbb{R}.
\]
\end{definition}

Let us test both notions on $X = \Fib$ with $T = S$ the shift map.

\begin{proposition}[Fixed Points of the Shift Map]
\label{prop:fixed_points_shift}
The only fixed point of the shift map $S$ is the zero sequence $\mathcal{F}_{0,0}$.
\end{proposition}

\begin{proof}
Suppose $\mathcal{A} = (a_0, a_1, a_2, \dots)$ satisfies $S(\mathcal{A}) = \mathcal{A}$.
Comparing the two sequences term by term gives
\[
a_0 = a_1, \quad a_1 = a_2, \quad a_2 = a_3, \quad \dots
\]
and therefore $a_n = a_0$ for all $n$. Such a constant sequence is a Fibonacci
sequence if and only if $a_0 = a_0 + a_0$, which forces $a_0 = 0$. Hence
$\mathcal{A} = (0,0,0,\dots) = \mathcal{F}_{0,0}$.
\end{proof}

So, the fixed points of $S$ are not interesting. Let us weaken the requirement and ask
for eigenvectors instead. This is exactly where the geometric sequences appear.

\begin{proposition}[Eigenvectors of the Shift Map]
\label{prop:eigenvectors_shift}
A sequence $\mathcal{A} \ne \mathcal{F}_{0,0}$ is an eigenvector of the shift map $S$
if and only if $\mathcal{A}$ is a geometric sequence
$\mathcal{A} = (a_0, a_0 \lambda, a_0 \lambda^2, \dots)$ with $a_0 \ne 0$ whose ratio
$\lambda$ satisfies $\lambda^2 = \lambda + 1$. Consequently
$\lambda \in \{\varphi, \psi\}$, while $a_0 \ne 0$ may be chosen freely.
\end{proposition}

\begin{proof}
Assume $S(\mathcal{A}) = \lambda \cdot \mathcal{A}$ for some $\lambda \in \mathbb{R}$.
Written out, the left-hand side is $(a_1, a_2, \dots, a_{n+1}, \dots)$ and the
right-hand side is $(\lambda a_0, \lambda a_1, \dots, \lambda a_n, \dots)$. Comparing
entries yields
\[
a_n = a_{n-1} \cdot \lambda \quad \text{for all } n \ge 1,
\]
and therefore
\[
\mathcal{A} = (a_0, \, a_0 \lambda, \, a_0 \lambda^2, \, a_0 \lambda^3, \, \dots, \, a_0 \lambda^n, \, \dots),
\]
which is precisely a geometric sequence. Note that $a_0 \ne 0$, since otherwise every
term would vanish and $\mathcal{A}$ would be $\mathcal{F}_{0,0}$, which we excluded.

It remains to determine which ratios $\lambda$ are admissible. Since $\mathcal{A}$ must
in addition be a Fibonacci sequence, the relation $a_2 = a_1 + a_0$ has to hold, that is
\[
\underbrace{a_0 \cdot \lambda^2}_{=\, a_2} = \underbrace{a_0 \cdot \lambda}_{=\, a_1} + \, a_0 .
\]
Dividing by $a_0 \ne 0$ gives $\lambda^2 = \lambda + 1$, the very same quadratic
equation we met in \cref{lem:geometric_fibonacci}. Its roots are $\varphi$ and $\psi$,
and $a_0$ is not constrained at all. Conversely, every geometric sequence with ratio
$\varphi$ or $\psi$ satisfies $S(\mathcal{A}) = \lambda \cdot \mathcal{A}$ by the same
computation read backwards.
\end{proof}

\begin{conclusion}
Taking $a_0 = 1$, the two geometric sequences
\[
\mathcal{F}_{1, \varphi} = (1, \varphi, \varphi^2, \dots, \varphi^n, \dots)
\qquad \text{and} \qquad
\mathcal{F}_{1, \psi} = (1, \psi, \psi^2, \dots, \psi^n, \dots)
\]
are Fibonacci sequences \textit{and} they are eigenvectors of the shift operator $S$.
The geometric sequences were therefore never a guess: they are forced upon us the
moment we ask which sequences the shift map merely rescales.
\end{conclusion}

It is worth restating this in the picture of \cref{fig:fib_plane}. Under the
identification $\mathcal{F}_{a_0,a_1} \leftrightarrow (a_0, a_1)$, shifting a sequence
discards $a_0$ and promotes $a_1$ and $a_2 = a_0 + a_1$ to the new initial pair, so
\begin{equation} \label{eq:shift_in_coordinates}
S\bigl( \mathcal{F}_{a_0, a_1} \bigr) = \mathcal{F}_{a_1, \, a_0 + a_1},
\qquad \text{that is} \qquad
\begin{pmatrix} a_0 \\ a_1 \end{pmatrix} \longmapsto
\begin{pmatrix} 0 & 1 \\ 1 & 1 \end{pmatrix}
\begin{pmatrix} a_0 \\ a_1 \end{pmatrix}.
\end{equation}
The shift map thus acts on the plane $\Fib$ through a single fixed matrix, and
\cref{prop:eigenvectors_shift} identifies the two lines that this matrix maps onto
themselves, as drawn in \cref{fig:shift_eigenlines}.

\begin{figure}[ht!]
  \centering
  \begin{tikzpicture}[scale=1.0, >=Stealth]
    %, the two invariant lines, labelled at their empty lower ends ---
    \draw[ThemeAccent!70, thick] (-1.20,-1.94) -- (1.88,3.04);
    \node[ThemeAccent, below right, font=\scriptsize] at (-1.22,-1.92) {$\mathbb{R} \cdot \mathcal{F}_{1,\varphi}$};
    \draw[ThemeGreen!70, thick] (-2.00,1.24) -- (2.90,-1.79);
    \node[ThemeGreen, below left, font=\scriptsize] at (2.92,-1.77) {$\mathbb{R} \cdot \mathcal{F}_{1,\psi}$};
    %, axes ---
    \draw[->, gray!70] (-2.20,0) -- (3.25,0) node[right, black, font=\small] {$a_0$};
    \draw[->, gray!70] (0,-2.15) -- (0,3.30) node[above, black, font=\small] {$a_1$};
    %, eigenvector for phi and its image (stretched by phi) ---
    \draw[->, ThemeAccent, very thick] (0,0) -- (1,1.618);
    \node[text=ThemeAccent, fill=white, inner sep=1pt, right, font=\scriptsize] at (1.10,1.50) {$\mathcal{F}_{1,\varphi}$};
    \draw[->, ThemeAccent, very thick, dashed] (0,0) -- (1.618,2.618);
    \node[ThemeAccent, right, font=\scriptsize] at (1.72,2.62) {$S(\mathcal{F}_{1,\varphi})$};
    %, eigenvector for psi and its image (flipped and contracted by psi) ---
    \draw[->, ThemeGreen, very thick] (0,0) -- (1,-0.618);
    \node[text=ThemeGreen, fill=white, inner sep=1pt, below right, font=\scriptsize] at (1.00,-0.82) {$\mathcal{F}_{1,\psi}$};
    \draw[->, ThemeGreen, very thick, dashed] (0,0) -- (-0.618,0.382);
    \node[text=ThemeGreen, fill=white, inner sep=1pt, above left, font=\scriptsize] at (-0.60,0.58) {$S(\mathcal{F}_{1,\psi})$};
    %, a generic sequence, whose direction is NOT preserved ---
    \draw[->, ThemePurple, thick] (0,0) -- (2.6,-0.5);
    \node[ThemePurple, below, font=\scriptsize] at (2.62,-0.58) {$\mathcal{A}$};
    \draw[->, ThemePurple, thick, dashed] (0,0) -- (-0.5,2.1);
    \node[ThemePurple, above left, font=\scriptsize] at (-0.48,2.12) {$S(\mathcal{A})$};
    \fill (0,0) circle (1.4pt);
  \end{tikzpicture}
  \caption{The shift map $S$ acting on the plane $\Fib$ via
  \eqref{eq:shift_in_coordinates}; images are drawn dashed. A generic sequence
  $\mathcal{A}$ (purple) is thrown off its own direction. Only the two coloured lines
  are preserved: on $\mathbb{R} \cdot \mathcal{F}_{1,\varphi}$ the map stretches by
  $\varphi > 1$, while on $\mathbb{R} \cdot \mathcal{F}_{1,\psi}$ it reverses the
  direction and contracts by
  $|\psi| < 1$. These two lines are exactly the geometric Fibonacci sequences of
  \cref{prop:eigenvectors_shift}.}
  \label{fig:shift_eigenlines}
\end{figure}

\begin{remark}
The contraction along the second line is the geometric reason behind the asymptotics
established below. Writing an arbitrary Fibonacci sequence in terms of the two
eigen-directions, the $\psi$-component is crushed towards $0$ under repeated shifting
while the $\varphi$-component grows. Whatever sequence one starts from, its direction
in the plane $\Fib$ converges to the line $\mathbb{R} \cdot \mathcal{F}_{1,\varphi}$.
\end{remark}

\subsection{The Upshot: Asymptotics and Matrix Form}

\begin{ainote}
This subsection has no counterpart in Prof.\ Biran's notes for Lecture 1A, which end
with the exercises reproduced below. The closest-integer estimate, the matrix form and
Cassini's identity are supplementary material added in this transcript. They are
included because the matrix $M$ reappears in the chapters on eigenvalues, but they
should not be mistaken for lecture content.
\end{ainote}

The derivation above yields two profound observations regarding the growth and representation of the Fibonacci sequence.

\begin{importantremark}
Clearly, as $n$ grows, the term involving $\psi$ becomes negligible. Since $|\psi| = |\frac{1-\sqrt{5}}{2}| \approx 0.618 < 1$, it follows that $\psi^n \to 0$ as $n \to \infty$. Therefore, for large $n$, we have the approximation $a_n \approx \frac{1}{\sqrt{5}} \varphi^n$. In fact, because $\left| \frac{\psi^n}{\sqrt{5}} \right| < \frac{1}{2}$ for all $n \ge 0$, we arrive at the following precise conclusion:
the number $a_n$ is the \textit{closest integer} to
\[
\frac{1}{\sqrt{5}} \cdot \left(\frac{1+\sqrt{5}}{2}\right)^n.
\]
Furthermore, this implies that the ratio of consecutive terms satisfies:
\begin{equation}
\lim_{n \to \infty} \frac{a_{n+1}}{a_n} = \varphi
\end{equation}
\end{importantremark}

\begin{figure}[ht!]
  \centering
  \begin{tikzpicture}[x=0.46cm, y=2.5cm, >=Stealth]
    %, limiting value phi ---
    \draw[ThemeAccent, dashed, thick] (0.4,1.61803) -- (12.1,1.61803);
    \node[ThemeAccent, right, font=\small] at (12.1,1.61803) {$\varphi$};
    %, axes ---
    \draw[->, gray!70] (0.4,0.95) -- (12.3,0.95) node[below, black, font=\small] {$n$};
    \draw[->, gray!70] (0.4,0.95) -- (0.4,2.34) node[left, black, font=\small] {$\dfrac{a_{n+1}}{a_n}$};
    %, y ticks ---
    \foreach \y in {1,1.5,2}
      \draw[gray!70] (0.32,\y) -- (0.48,\y) node[left=2pt, black, font=\scriptsize] {$\y$};
    %, x ticks ---
    \foreach \n in {2,4,6,8,10}
      \draw[gray!70] (\n,0.90) -- (\n,1.00) node[below=3pt, black, font=\scriptsize] {$\n$};
    %, the ratios a_{n+1}/a_n ---
    \draw[ThemePurple, thin]
      (1,1) -- (2,2) -- (3,1.5) -- (4,1.66667) -- (5,1.6) -- (6,1.625)
      -- (7,1.61538) -- (8,1.61905) -- (9,1.61765) -- (10,1.61818) -- (11,1.61798);
    \foreach \p in {(1,1),(2,2),(3,1.5),(4,1.66667),(5,1.6),(6,1.625),%
                    (7,1.61538),(8,1.61905),(9,1.61765),(10,1.61818),(11,1.61798)}
      \fill[ThemePurple] \p circle (1.6pt);
  \end{tikzpicture}
  \caption{The quotients $a_{n+1}/a_n$ of the classical Fibonacci sequence
  $\mathcal{F}_{0,1}$. They oscillate around the golden ratio $\varphi$, alternately
  overshooting and undershooting it, because the error term $\psi^n$ changes sign at
  every step while its magnitude decays geometrically.}
  \label{fig:fib_ratio_to_phi}
\end{figure}

\begin{exercise}[Asymptotic Ratio of Consecutive Fibonacci Numbers]
\label{exc:golden_ratio_part2}
Let $\mathcal{F}_{0,1} = (F_0, F_1, F_2, \dots) \in \Fib$ be the classical Fibonacci sequence. Assuming that the limit
\[
L := \lim_{n \to \infty} \frac{F_n}{F_{n-1}}
\]
exists and is non-zero, compute its exact value.
\exinfo{This exercise is Problem 3 of Problem Sheet 1, Linear Algebra I, Autumn Semester 2025. Taken from Dr.~Menny Akka Ginosar's notes on Fibonacci sequences.}
\exsol{sol:golden_ratio_part2}
\end{exercise}

Finally, we can represent the recurrence relation using the language of matrices. Defining a vector $v_n := \binom{a_{n+1}}{a_n}$, the relation $a_{n+1} = a_n + a_{n-1}$ can be written as:
\begin{equation}
\begin{pmatrix} a_{n+1} \\ a_n \end{pmatrix} = \begin{pmatrix} 1 & 1 \\ 1 & 0 \end{pmatrix} \begin{pmatrix} a_n \\ a_{n-1} \end{pmatrix} \implies v_n = \begin{pmatrix} 1 & 1 \\ 1 & 0 \end{pmatrix}^n \begin{pmatrix} a_1 \\ a_0 \end{pmatrix}
\end{equation}
This formulation shows that the behavior of the sequence is entirely governed by the powers of the matrix $M := \left(\begin{smallmatrix} 1 & 1 \\ 1 & 0 \end{smallmatrix}\right)$. It is the shift map of \eqref{eq:shift_in_coordinates} written in the reversed coordinates $\binom{a_{n+1}}{a_n}$, so by \cref{prop:eigenvectors_shift} its eigenvalues are again $\varphi$ and $\psi$.

\begin{exercise}[Further Problems on Fibonacci Sequences and the Golden Ratio]
\label{exc:ch1_list}
\begin{enumerate}
    \item \textbf{General Formula:} Find an explicit formula for the general Fibonacci sequence $\mathcal{F}_{a_0, a_1}$ that depends only on the initial values $a_0, a_1$ and the index $n$.
    \item \textbf{The Shift Operator:} Let $S: \Fib \to \Fib$ be the shift map defined by $S(a_0, a_1, \dots) = (a_1, a_2, \dots)$. Formally, verify that $S$ is a \textit{linear map}.
    \item \textbf{Golden Ratio Continued Fraction:} Consider $\varphi = \frac{1+\sqrt{5}}{2}$.
        \[
        \left.
        \begin{aligned}
        \varphi
        = 1+\cfrac{1}{\,1+\cfrac{1}{\,1+\cfrac{1}{\,1+\cfrac{1}{\,1+\cdots}}}}
        \end{aligned}
        \qquad
        \right\}
        \qquad
        \left.
        \begin{aligned}
        &\text{what is actually}\\
        &\text{the meaning of}\\
        &\text{this?}
        \end{aligned}
        \right.
        \]
    \item \textbf{Non-existence of Arithmetic Fibonacci Sequences:} Complete the details to show that no Fibonacci sequence $\mathcal{F}_{a_0, a_1}$ (except for the trivial zero sequence) can be an arithmetic sequence.
    \item \textbf{Convergence:} For the classical Fibonacci sequence $\mathcal{F}_{0,1} = (a_0, a_1, a_2, \dots)$, calculate the limit
    \[
        \lim_{n \to \infty} \frac{a_n}{a_{n-1}},
    \]
    whose numerical behaviour is plotted in \cref{fig:fib_ratio_to_phi}.
    \item \textbf{Cassini's Identity:} Show that $a_{n+1} a_{n-1} - a_n^2 = (-1)^n$ for all $n \ge 1$. \exhint[Hint (Prof. Biran)]{Use the determinant of the matrix $M^n$.}
\end{enumerate}
\exinfo{This exercise is taken from Prof.~Biran's handwritten lecture notes.}
\exsol{sol:shift_is_linear_b}
\end{exercise}

% Creator: Opus 5
\subsection{Solutions to the Exercises}
\label{sec:solutions_fibonacci}

\begin{exercisesolution}[Proof of \cref{exc:shift_is_linear}, also \cref{exc:ch1_list} \textbf{(b)}]
\label{sol:shift_is_linear_b}
Let $\mathcal{A} = (a_0, a_1, a_2, \dots)$ and $\mathcal{B} = (b_0, b_1, b_2, \dots)$ be in $\Fib$ and let $\alpha \in \mathbb{R}$. Both operations on $\Fib$ are term-wise, and the shift acts on each term separately, so the two axioms of \cref{def:linear_map_fib} are read off directly:
\[
    S(\mathcal{A} + \mathcal{B}) = S(a_0 + b_0, \, a_1 + b_1, \, \dots) = (a_1 + b_1, \, a_2 + b_2, \, \dots) = S(\mathcal{A}) + S(\mathcal{B}),
\]
\[
    S(\alpha \mathcal{A}) = S(\alpha a_0, \, \alpha a_1, \, \dots) = (\alpha a_1, \, \alpha a_2, \, \dots) = \alpha \, S(\mathcal{A}).
\]
Both computations use \cref{prop:fib_closure} implicitly, in that the sequences $\mathcal{A} + \mathcal{B}$ and $\alpha \mathcal{A}$ are again in $\Fib$, so that $S$ may be applied to them at all.
\end{exercisesolution}

\begin{exercisesolution}[Solution of \cref{exc:ch1_list} \textbf{(a)}]
\label{sol:ch1_list_a}
We repeat the derivation of \cref{thm:binet_formula} with $\mathcal{F}_{0,1}$ replaced by an arbitrary $\mathcal{F}_{a_0, a_1}$. By \cref{cor:fib_linear_combination} we look for $\alpha, \beta \in \mathbb{R}$ with $\alpha \cdot \mathcal{F}_{1,\varphi} + \beta \cdot \mathcal{F}_{1,\psi} = \mathcal{F}_{a_0, a_1}$, and comparing the terms with $n = 0$ and $n = 1$ gives
\[
\begin{cases}
    \alpha + \beta = a_0, \
    \alpha \varphi + \beta \psi = a_1.
\end{cases}
\]
Substituting $\beta = a_0 - \alpha$ into the second equation gives $\alpha(\varphi - \psi) = a_1 - a_0 \psi$, and since $\varphi - \psi = \sqrt{5} \neq 0$,
\[
    \alpha = \frac{a_1 - a_0 \psi}{\sqrt{5}}, \qquad \beta = a_0 - \alpha = \frac{a_0 \varphi - a_1}{\sqrt{5}}.
\]
Reading off the $n$-th term, the general Fibonacci sequence $\mathcal{F}_{a_0, a_1} = (\alpha_0, \alpha_1, \alpha_2, \dots)$ satisfies
\[
    \alpha_n = \frac{1}{\sqrt{5}} \left( (a_1 - a_0 \psi) \, \varphi^n - (a_1 - a_0 \varphi) \, \psi^n \right) \qquad \text{for all } n \ge 0.
\]
Two checks. For $a_0 = 0$ and $a_1 = 1$ this collapses to $\frac{1}{\sqrt{5}}(\varphi^n - \psi^n)$, which is Binet's formula. For $a_0 = 1$ and $a_1 = 0$ it gives $\frac{1}{\sqrt{5}}(\varphi \psi^n - \psi \varphi^n)$, whose values at $n = 0$ and $n = 1$ are $\frac{\varphi - \psi}{\sqrt{5}} = 1$ and $\frac{\varphi\psi - \psi\varphi}{\sqrt{5}} = 0$, as they must be.
\end{exercisesolution}

\begin{exercisesolution}[Solution of \cref{exc:ch1_list} \textbf{(c)}]
\label{sol:ch1_list_c}
The displayed expression is not a finite computation, so it has to be given a meaning, and the standard one is as a limit. Define a sequence by
\[
    x_1 := 1, \qquad x_{n+1} := 1 + \frac{1}{x_n},
\]
so that $x_n$ is the continued fraction truncated after $n$ layers, and read the infinite fraction as $\lim_{n \to \infty} x_n$, provided that limit exists.

Two things then have to be said. First, \emph{if} the limit $x$ exists, it is $\varphi$: every $x_n \ge 1$, so $x \ge 1 > 0$, and passing to the limit in the recursion gives $x = 1 + \frac{1}{x}$, that is, $x^2 = x + 1$. That is exactly the equation of \cref{lem:geometric_fibonacci}, whose roots are $\varphi$ and $\psi$; since $\psi < 0$ and $x > 0$, we get $x = \varphi$.

Second, the limit does exist, and the reason is this chapter's own subject. Starting from $x_1 = 1 = a_2/a_1$, an induction gives
\[
    x_n = \frac{a_{n+1}}{a_n},
\]
because $1 + \frac{a_n}{a_{n+1}} = \frac{a_{n+1} + a_n}{a_{n+1}} = \frac{a_{n+2}}{a_{n+1}}$ by the recurrence. So, the convergents of the continued fraction are precisely the ratios of consecutive Fibonacci numbers, and part \textbf{(e)} below shows that these converge to $\varphi$.

The meaning of the picture, then, is that the golden ratio is the number whose continued fraction expansion consists of nothing but $1$'s, and that its convergents are the Fibonacci quotients plotted in \cref{fig:fib_ratio_to_phi}.
\end{exercisesolution}

\begin{exercisesolution}[Solution of \cref{exc:ch1_list} \textbf{(d)}]
\label{sol:ch1_list_d}
Let $\mathcal{F}_{a_0, a_1} = (\alpha_0, \alpha_1, \alpha_2, \dots)$ be a Fibonacci sequence which is also arithmetic, say with common difference $d$, so that $\alpha_n - \alpha_{n-1} = d$ for all $n \ge 1$. For $n \ge 2$ the Fibonacci recurrence gives $\alpha_n - \alpha_{n-1} = \alpha_{n-2}$, and comparing the two expressions yields
\[
    \alpha_{n-2} = d \qquad \text{for all } n \ge 2,
\]
that is, $\alpha_m = d$ for every $m \ge 0$: the sequence is constant. Now impose the recurrence once on this constant sequence: $d = \alpha_2 = \alpha_1 + \alpha_0 = d + d$, so $d = 0$ and $\mathcal{F}_{a_0,a_1} = \mathcal{F}_{0,0}$.

Conversely, $\mathcal{F}_{0,0}$ is arithmetic, with $d = 0$. So, the zero sequence is the only arithmetic Fibonacci sequence, which is what the text asserts before \cref{lem:geometric_fibonacci}. Note that no growth estimate entered: the argument is finished by the single equation $d = d + d$, and it applies to sequences with negative terms just as well.
\end{exercisesolution}

\begin{exercisesolution}[Solution of \cref{exc:ch1_list} \textbf{(e)}]
\label{sol:ch1_list_e}
By \cref{thm:binet_formula}, for $n \ge 1$,
\[
    \frac{a_n}{a_{n-1}} = \frac{\varphi^n - \psi^n}{\varphi^{n-1} - \psi^{n-1}} = \varphi \cdot \frac{1 - (\psi/\varphi)^n}{1 - (\psi/\varphi)^{n-1}} \cdot \frac{1}{1},
\]
where the second expression is obtained by dividing numerator and denominator by $\varphi^n$ and $\varphi^{n-1}$ respectively. Now $\left|\psi/\varphi\right| = \frac{|1 - \sqrt{5}|}{1 + \sqrt{5}} = \frac{\sqrt{5}-1}{\sqrt{5}+1} < 1$, so $(\psi/\varphi)^n \to 0$, and both the numerator and the denominator of the middle fraction tend to $1$. Note also that $a_{n-1} \neq 0$ for $n \ge 2$, so the quotient is defined. Consequently
\[
    \lim_{n \to \infty} \frac{a_n}{a_{n-1}} = \varphi = \frac{1 + \sqrt{5}}{2}.
\]
The oscillation visible in \cref{fig:fib_ratio_to_phi} is the sign of $(\psi/\varphi)^n$, which alternates, while its magnitude decays geometrically.
\end{exercisesolution}

\begin{exercisesolution}[Solution of \cref{exc:ch1_list} \textbf{(f)}]
\label{sol:ch1_list_f}
Write $M := \left(\begin{smallmatrix} 1 & 1 \\ 1 & 0 \end{smallmatrix}\right)$ as in the text. We first claim that
\[
    M^n = \begin{pmatrix} a_{n+1} & a_n \\ a_n & a_{n-1} \end{pmatrix} \qquad \text{for all } n \ge 1,
\]
where $(a_n)$ is the classical Fibonacci sequence $\mathcal{F}_{0,1}$. For $n = 1$ the right-hand side is $\left(\begin{smallmatrix} a_2 & a_1 \\ a_1 & a_0 \end{smallmatrix}\right) = \left(\begin{smallmatrix} 1 & 1 \\ 1 & 0 \end{smallmatrix}\right) = M$. Assuming the claim for $n$,
\begin{align*}
    M^{n+1} = M^n \cdot M
    &= \begin{pmatrix} a_{n+1} & a_n \\ a_n & a_{n-1} \end{pmatrix} \begin{pmatrix} 1 & 1 \\ 1 & 0 \end{pmatrix} \\
    &= \begin{pmatrix} a_{n+1} + a_n & a_{n+1} \\ a_n + a_{n-1} & a_n \end{pmatrix}
     = \begin{pmatrix} a_{n+2} & a_{n+1} \\ a_{n+1} & a_n \end{pmatrix},
\end{align*}
using the recurrence twice, which is the claim for $n+1$.

Now take determinants. The determinant is multiplicative, so $\det(M^n) = (\det M)^n$, and $\det M = 1 \cdot 0 - 1 \cdot 1 = -1$. Reading the determinant off the claim instead gives $a_{n+1} a_{n-1} - a_n^2$. Equating the two,
\[
    a_{n+1} a_{n-1} - a_n^2 = (-1)^n \qquad \text{for all } n \ge 1,
\]
which is Cassini's identity. For instance, $n = 6$ gives $a_7 a_5 - a_6^2 = 13 \cdot 5 - 8^2 = 65 - 64 = 1 = (-1)^6$.
\end{exercisesolution}

\begin{ainote}
This is the only chapter that answers a question rather than building a theory,
and everything in it is one move: replace a question about the single number
$a_n$ by a question about the whole space $\Fib$ of sequences like it. That space
is closed under addition and scaling (\cref{prop:fib_closure}), so a formula for
two convenient sequences yields one for all of them. The convenient ones are the
geometric sequences, and \cref{prop:eigenvectors_shift} says why they were not a
lucky guess: they are exactly the sequences the shift map rescales. Every later
chapter repeats the move with \qt{vector space}, \qt{basis}, \qt{linear map} and
\qt{eigenvector} in place of the ad-hoc words used here.

One caution about the text. The argument ruling out arithmetic Fibonacci
sequences appeals to the difference $a_{n-2}$ \qt{growing}, which is not what
makes it work: the displayed identities say the sequence is \emph{constant}, and
a constant sequence obeys the recurrence only when $d = d + d$. No growth
estimate enters, and the conclusion covers sequences with negative terms. This is
\cref{exc:ch1_list} \textbf{(d)}; the details are in the solutions above.
\end{ainote}

% Solution by Gemini 3.7 Flash (High)
\begin{exercisesolution}[Solution to \cref{exc:fib_dimension_two_samples}]
\label{sol:fib_dimension_two_samples}
Let $F = (\alpha_n)_{n=0}^\infty$ and $G = (\beta_n)_{n=0}^\infty$ be two sequences in $\Fib$. Consider the difference sequence $D := F - G = (\delta_n)_{n=0}^\infty$, where $\delta_n := \alpha_n - \beta_n$ for all $n \ge 0$. 

By \cref{prop:fib_closure}, $\Fib$ is closed under linear combinations, so $D \in \Fib$. Therefore, by \cref{thm:binet_formula}, there exist unique real constants $c_1, c_2 \in \mathbb{R}$ such that
\[
\delta_n = c_1 \varphi^n + c_2 \psi^n \quad \text{for all } n \ge 0,
\]
where $\varphi = \frac{1+\sqrt{5}}{2}$ and $\psi = \frac{1-\sqrt{5}}{2}$ are the two roots of $x^2 - x - 1 = 0$.

Suppose we check equality at two distinct indices $i < j$. The condition $\alpha_i = \beta_i$ and $\alpha_j = \beta_j$ is equivalent to $\delta_i = 0$ and $\delta_j = 0$, which yields the linear system of equations in the unknowns $(c_1, c_2)$:
\[
\begin{cases}
c_1 \varphi^i + c_2 \psi^i = 0, \\
c_1 \varphi^j + c_2 \psi^j = 0.
\end{cases}
\]
In matrix form, this reads:
\[
\begin{pmatrix}
\varphi^i & \psi^i \\
\varphi^j & \psi^j
\end{pmatrix}
\begin{pmatrix}
c_1 \\
c_2
\end{pmatrix}
=
\begin{pmatrix}
0 \\
0
\end{pmatrix}.
\]
The determinant of this $2 \times 2$ coefficient matrix is:
\begin{align*}
\det \begin{pmatrix} \varphi^i & \psi^i \\ \varphi^j & \psi^j \end{pmatrix} &= \varphi^i \psi^j - \varphi^j \psi^i \\
&= \varphi^i \psi^i \left(\psi^{j-i} - \varphi^{j-i}\right).
\end{align*}
Since $\varphi \cdot \psi = \frac{1-5}{4} = -1 \ne 0$, we have $\varphi^i \psi^i = (-1)^i \ne 0$. Furthermore, because $\varphi > 1$ and $-1 < \psi < 0$, we have $|\varphi| > |\psi|$, and thus for any integer $k = j - i > 0$, $\varphi^k \ne \psi^k$. Hence,
\[
\psi^{j-i} - \varphi^{j-i} \ne 0,
\]
which proves that the determinant is non-zero.

A $2 \times 2$ homogeneous linear system with a non-zero determinant has only the trivial solution $c_1 = c_2 = 0$. Consequently, $\delta_n = 0$ for all $n \ge 0$, which implies $D = \mathcal{F}_{0,0}$ and therefore $F = G$. Thus, sampling two distinct indices is always sufficient to determine equality of Fibonacci sequences.

\begin{ainote}
\textbf{Alternative Elementary Proof via Recurrence Shifting:}
An alternative perspective avoids the geometric eigenbasis and works directly with the recurrence relation. Let $\delta_i = 0$. Using the recurrence $\delta_{m+2} = \delta_{m+1} + \delta_m$, we can compute both forwards and backwards from index $i$:
\begin{align*}
\text{Forward: } \delta_{i+1} &= \delta_{i+1}, \quad \delta_{i+2} = \delta_{i+1}, \quad \delta_{i+3} = 2\delta_{i+1}, \quad \dots \quad \delta_{i+k} = F_k \delta_{i+1}, \\
\text{Backward: } \delta_{i-1} &= \delta_{i+1} - \delta_i = \delta_{i+1}, \quad \delta_{i-2} = -\delta_{i+1}, \quad \dots \quad \delta_{i-\ell} = (-1)^{\ell+1} F_\ell \delta_{i+1},
\end{align*}
where $(F_k)$ is the standard Fibonacci sequence $\mathcal{F}_{0,1} = (0, 1, 1, 2, 3, \dots)$.
Since $j > i$, the second condition gives $0 = \delta_j = \delta_{i+(j-i)} = F_{j-i} \delta_{i+1}$. Because $j - i \ge 1$, we have $F_{j-i} > 0$, forcing $\delta_{i+1} = 0$. Since $\delta_i = 0$ and $\delta_{i+1} = 0$, every term $\delta_n = 0$, so $F = G$.
\end{ainote}
\end{exercisesolution}

% Solution by Gemini 3.7 Flash (High)
\begin{exercisesolution}[Solution to \cref{exc:golden_ratio_part1}]
\label{sol:golden_ratio_part1}
\begin{enumerate}[label=\textbf{(\alph*)}]
    \item We seek the explicit formula for $\mathcal{F}_{1,0} = (b_n)_{n=0}^\infty$. Since $\mathcal{F}_{1,0} \in \Fib$, we can express it as a linear combination of the two geometric basis sequences $\mathcal{F}_{1,\varphi} = (\varphi^n)$ and $\mathcal{F}_{1,\psi} = (\psi^n)$:
    \[
    b_n = c_1 \varphi^n + c_2 \psi^n.
    \]
    Using the initial conditions $b_0 = 1$ and $b_1 = 0$, we obtain:
    \[
    \begin{cases}
    c_1 + c_2 = 1, \\
    c_1 \varphi + c_2 \psi = 0.
    \end{cases}
    \]
    From the second equation, $c_2 = -\frac{\varphi}{\psi} c_1$. Substituting this into the first equation:
    \[
    c_1 \left(1 - \frac{\varphi}{\psi}\right) = 1 \implies c_1 \left(\frac{\psi - \varphi}{\psi}\right) = 1 \implies c_1 = \frac{\psi}{\psi - \varphi}.
    \]
    Since $\psi - \varphi = \frac{1-\sqrt{5}}{2} - \frac{1+\sqrt{5}}{2} = -\sqrt{5}$, we find
    \[
    c_1 = \frac{\psi}{-\sqrt{5}} = -\frac{\psi}{\sqrt{5}} = \frac{\sqrt{5}-1}{2\sqrt{5}}, \quad c_2 = 1 - c_1 = \frac{\varphi}{\sqrt{5}} = \frac{\sqrt{5}+1}{2\sqrt{5}}.
    \]
    Therefore, the $n$-th term of $\mathcal{F}_{1,0}$ is:
    \[
    b_n = -\frac{\psi}{\sqrt{5}} \varphi^n + \frac{\varphi}{\sqrt{5}} \psi^n = \frac{\varphi \psi (\psi^{n-1} - \varphi^{n-1})}{\sqrt{5}}.
    \]
    Using $\varphi \psi = -1$, this simplifies for all $n \ge 1$ to:
    \[
    b_n = \frac{\varphi^{n-1} - \psi^{n-1}}{\sqrt{5}} = a_{n-1},
    \]
    where $a_n = \frac{\varphi^n - \psi^n}{\sqrt{5}}$ is the classical Fibonacci sequence $\mathcal{F}_{0,1}$. Note that for $n = 0$, this yields $b_0 = 1$.

    \item Let $\mathcal{F}_{a,b} = (x_n)_{n=0}^\infty$ be the generalized Fibonacci sequence with initial values $x_0 = a$ and $x_1 = b$. By \cref{cor:fib_linear_combination},
    \[
    \mathcal{F}_{a,b} = a \mathcal{F}_{1,0} + b \mathcal{F}_{0,1}.
    \]
    Therefore, for every $n \ge 0$,
    \[
    x_n = a b_n + b a_n = a \left(\frac{-\psi \varphi^n + \varphi \psi^n}{\sqrt{5}}\right) + b \left(\frac{\varphi^n - \psi^n}{\sqrt{5}}\right) = \left(\frac{b - a\psi}{\sqrt{5}}\right) \varphi^n + \left(\frac{a\varphi - b}{\sqrt{5}}\right) \psi^n.
    \]
    Using $\psi = \frac{1-\sqrt{5}}{2}$ and $\varphi = \frac{1+\sqrt{5}}{2}$, we can also write this as:
    \[
    x_n = \frac{1}{\sqrt{5}} \left[ \left(b - a \frac{1-\sqrt{5}}{2}\right) \left(\frac{1+\sqrt{5}}{2}\right)^n + \left(a \frac{1+\sqrt{5}}{2} - b\right) \left(\frac{1-\sqrt{5}}{2}\right)^n \right].
    \]
\end{enumerate}
\end{exercisesolution}

% Solution by Gemini 3.7 Flash (High)
\begin{exercisesolution}[Solution to \cref{exc:golden_ratio_part2}]
\label{sol:golden_ratio_part2}
Let $F_n$ denote the $n$-th classical Fibonacci number ($F_0 = 0, F_1 = 1, F_2 = 1, F_3 = 2, \dots$). We are given that the limit
\[
L := \lim_{n \to \infty} \frac{F_n}{F_{n-1}}
\]
exists and is non-zero. Dividing the Fibonacci recurrence $F_n = F_{n-1} + F_{n-2}$ by $F_{n-1}$ (which is positive for $n \ge 2$) gives:
\[
\frac{F_n}{F_{n-1}} = 1 + \frac{F_{n-2}}{F_{n-1}} = 1 + \frac{1}{\frac{F_{n-1}}{F_{n-2}}}.
\]
Taking the limit as $n \to \infty$ on both sides and using the algebraic limit laws (since $L \ne 0$), we obtain:
\[
L = 1 + \frac{1}{L} \implies L^2 - L - 1 = 0.
\]
The solutions to this quadratic equation are $L = \frac{1 \pm \sqrt{5}}{2}$. Since $F_n > 0$ for all $n \ge 1$, the ratio $\frac{F_n}{F_{n-1}} > 0$ is strictly positive for all $n \ge 2$, so the limit must satisfy $L \ge 0$. This rules out the negative root $\frac{1-\sqrt{5}}{2}$.

Therefore, the exact value of the limit is the golden ratio:
\[
L = \varphi = \frac{1 + \sqrt{5}}{2} \approx 1.6180339887\dots
\]
\end{exercisesolution}



